\section{AMC 12 12A 2026}
        \begin{enumerate}
        	\item (\textit{AMC 12 12A 2026}) A sign in a parking lot states that the total price includes a $27\%$ tax. To the nearest whole number, what percentage of the amount a customer pays is this tax?


 $\textbf{(A)}\   21\%  \qquad 	\textbf{(B)}\   22\%  \qquad 	\textbf{(C)}\   27\%  \qquad 	\textbf{(D)}\   37\%  \qquad 	\textbf{(E)}\   73\%$


	\item (\textit{AMC 12 12A 2026}) Square $ABCD$ has side length $24$. Let $E$ be the midpoint of $\overline{BC}$, and let $F$ be the intersection of $\overline{AE}$ and $\overline{BD}$.


 
                    \begin{figure}[H]
                        \centering
                        \includegraphics[width=0.5\textwidth]{Problem/AMC_12_12A_2026_2_Problem_Figure_1.png}
                    \end{figure}
                    


 What is the area of $\triangle BEF$?


 $\textbf{(A)}\   36  \qquad 	\textbf{(B)}\   40  \qquad 	\textbf{(C)}\   45  \qquad 	\textbf{(D)}\   48  \qquad 	\textbf{(E)}\   52$


	\item (\textit{AMC 12 12A 2026}) Half of the internal angles of a polygon measure $150^{\circ}$ each, and the other half of its internal angles measure $190^{\circ}$ each. How many sides does the polygon have?


 $\textbf{(A)}\   36  \qquad 	\textbf{(B)}\   40  \qquad 	\textbf{(C)}\   60  \qquad 	\textbf{(D)}\   72  \qquad 	\textbf{(E)}\   180$


	\item (\textit{AMC 12 12A 2026}) Let $a_{1}, a_{2}, a_{3}$, $\dots$ be a sequence with $a_{1} = 94$, $a_{2} = 190$, and $a_{n}$ equal to the arithmetic mean of the previous $n-1$ terms for $n \ge 3$. What is the value of $a_{10}$?


 $\textbf{(A)}\   47  \qquad 	\textbf{(B)}\   95  \qquad 	\textbf{(C)}\   142  \qquad 	\textbf{(D)}\   284  \qquad 	\textbf{(E)}\   564$


	\item (\textit{AMC 12 12A 2026}) How many ordered triples of positive integers $(x,y,z)$ satisfy both of the following equations?
<mjx-container class="MathJax CtxtMenu\_Attached\_0" ctxtmenu\_counter="0" display="true" jax="CHTML" style="font-size: 117.4\%; position: relative;" tabindex="0"><mjx-math aria-hidden="true" class="MJX-TEX" display="true" style="margin-left: 0px; margin-right: 0px;"><mjx-mtable style="min-width: 7.305em;"><mjx-table><mjx-itable><mjx-mtr><mjx-mtd style="text-align: right; padding-right: 0px; padding-bottom: 0.15em;"><mjx-mi class="mjx-i"><mjx-c class="mjx-c1D465 TEX-I"></mjx-c></mjx-mi><mjx-mo class="mjx-n" space="3"><mjx-c class="mjx-c2B"></mjx-c></mjx-mo><mjx-mi class="mjx-i" space="3"><mjx-c class="mjx-c1D466 TEX-I"></mjx-c></mjx-mi><mjx-mo class="mjx-n" space="3"><mjx-c class="mjx-c2B"></mjx-c></mjx-mo><mjx-mi class="mjx-i" space="3"><mjx-c class="mjx-c1D467 TEX-I"></mjx-c></mjx-mi><mjx-tstrut></mjx-tstrut></mjx-mtd><mjx-mtd style="text-align: left; padding-left: 0px; padding-bottom: 0.15em;"><mjx-mi class="mjx-n"></mjx-mi><mjx-mo class="mjx-n" space="4"><mjx-c class="mjx-c3D"></mjx-c></mjx-mo><mjx-mn class="mjx-n" space="4"><mjx-c class="mjx-c33"></mjx-c><mjx-c class="mjx-c31"></mjx-c></mjx-mn><mjx-tstrut></mjx-tstrut></mjx-mtd></mjx-mtr><mjx-mtr><mjx-mtd style="text-align: right; padding-right: 0px; padding-top: 0.15em;"><mjx-mi class="mjx-i"><mjx-c class="mjx-c1D465 TEX-I"></mjx-c></mjx-mi><mjx-mo class="mjx-n" space="3"><mjx-c class="mjx-c2B"></mjx-c></mjx-mo><mjx-mn class="mjx-n" space="3"><mjx-c class="mjx-c33"></mjx-c></mjx-mn><mjx-mi class="mjx-i"><mjx-c class="mjx-c1D466 TEX-I"></mjx-c></mjx-mi><mjx-mo class="mjx-n" space="3"><mjx-c class="mjx-c2B"></mjx-c></mjx-mo><mjx-mn class="mjx-n" space="3"><mjx-c class="mjx-c33"></mjx-c></mjx-mn><mjx-mi class="mjx-i"><mjx-c class="mjx-c1D467 TEX-I"></mjx-c></mjx-mi><mjx-tstrut></mjx-tstrut></mjx-mtd><mjx-mtd style="text-align: left; padding-left: 0px; padding-top: 0.15em;"><mjx-mi class="mjx-n"></mjx-mi><mjx-mo class="mjx-n" space="4"><mjx-c class="mjx-c3D"></mjx-c></mjx-mo><mjx-mn class="mjx-n" space="4"><mjx-c class="mjx-c34"></mjx-c><mjx-c class="mjx-c31"></mjx-c></mjx-mn><mjx-tstrut></mjx-tstrut></mjx-mtd></mjx-mtr></mjx-itable></mjx-table></mjx-mtable></mjx-math><mjx-assistive-mml display="block" unselectable="on"><math display="block" xmlns="http://www.w3.org/1998/Math/MathML"><mtable columnalign="right left" columnspacing="0em" displaystyle="true" rowspacing="3pt"><mtr><mtd><mi>x</mi><mo>+</mo><mi>y</mi><mo>+</mo><mi>z</mi></mtd><mtd><mi></mi><mo>=</mo><mn>31</mn></mtd></mtr><mtr><mtd><mi>x</mi><mo>+</mo><mn>3</mn><mi>y</mi><mo>+</mo><mn>3</mn><mi>z</mi></mtd><mtd><mi></mi><mo>=</mo><mn>41</mn></mtd></mtr></mtable></math></mjx-assistive-mml></mjx-container>
$\textbf{(A)}\   4  \qquad 	\textbf{(B)}\   5  \qquad 	\textbf{(C)}\   6  \qquad 	\textbf{(D)}\   7  \qquad 	\textbf{(E)}\   8$


	\item (\textit{AMC 12 12A 2026}) For what positive integer $n$ is the mean of $1$, $4$, $9$, $16$, $\dots$, $n^{2}$ exactly $32\%$ greater than its median?


 $\textbf{(A)}\   19  \qquad 	\textbf{(B)}\   39  \qquad 	\textbf{(C)}\   49  \qquad 	\textbf{(D)}\   89  \qquad 	\textbf{(E)}\   169$


	\item (\textit{AMC 12 12A 2026}) What is the value of $(1+i)^{20} + (1-i)^{20}$, where $i = \sqrt{-1}$?


 $\textbf{(A)}\   -2048  \qquad 	\textbf{(B)}\   -1024  \qquad 	\textbf{(C)}\   0  \qquad 	\textbf{(D)}\   1024  \qquad 	\textbf{(E)}\   2048$


	\item (\textit{AMC 12 12A 2026}) At Euclid High School, 25 students registered to take the AMC 10 and 75 students registered to take the AMC 12. (No student registered to take both.) However, not all 100 students showed up on test day. To compute the percentage of those students who actually took one of the exams, Euclid High School performed the following incorrect two-step process:


 
            \begin{itemize}
        \item Step 1: The school computed both the percentage of AMC 10 registrants who actually took that exam and the percentage of AMC 12 registrants who actually took that exam.
            \end{itemize}
            
            \begin{itemize}
        \item Step 2: The school took the average of the previous two numbers to obtain a final percentage.
            \end{itemize}
            Suppose that Euclid High School's process yielded an answer of $60\%$. What is the minimum number of students who might have actually taken one of the exams?


 $\textbf{(A)}\   30  \qquad 	\textbf{(B)}\   35  \qquad 	\textbf{(C)}\   40  \qquad 	\textbf{(D)}\   50  \qquad 	\textbf{(E)}\   60$


	\item (\textit{AMC 12 12A 2026}) How many ordered pairs $(x,y)$ of real numbers satisfy the following system of equations?
<mjx-container class="MathJax CtxtMenu\_Attached\_0" ctxtmenu\_counter="1" display="true" jax="CHTML" style="font-size: 117.4\%; position: relative;" tabindex="0"><mjx-math aria-hidden="true" class="MJX-TEX" display="true" style="margin-left: 0px; margin-right: 0px;"><mjx-mtable style="min-width: 6.109em;"><mjx-table><mjx-itable><mjx-mtr><mjx-mtd style="text-align: right; padding-right: 0px; padding-bottom: 0.15em;"><mjx-mo class="mjx-n"><mjx-c class="mjx-c7C"></mjx-c></mjx-mo><mjx-msup><mjx-mi class="mjx-i"><mjx-c class="mjx-c1D465 TEX-I"></mjx-c></mjx-mi><mjx-script style="vertical-align: 0.413em;"><mjx-texatom size="s" texclass="ORD"><mjx-mn class="mjx-n"><mjx-c class="mjx-c32"></mjx-c></mjx-mn></mjx-texatom></mjx-script></mjx-msup><mjx-mo class="mjx-n" space="3"><mjx-c class="mjx-c2212"></mjx-c></mjx-mo><mjx-msup space="3"><mjx-mi class="mjx-i"><mjx-c class="mjx-c1D466 TEX-I"></mjx-c></mjx-mi><mjx-script style="vertical-align: 0.413em;"><mjx-texatom size="s" texclass="ORD"><mjx-mn class="mjx-n"><mjx-c class="mjx-c32"></mjx-c></mjx-mn></mjx-texatom></mjx-script></mjx-msup><mjx-mo class="mjx-n"><mjx-c class="mjx-c7C"></mjx-c></mjx-mo><mjx-tstrut></mjx-tstrut></mjx-mtd><mjx-mtd style="text-align: left; padding-left: 0px; padding-bottom: 0.15em;"><mjx-mi class="mjx-n"></mjx-mi><mjx-mo class="mjx-n" space="4"><mjx-c class="mjx-c3D"></mjx-c></mjx-mo><mjx-mn class="mjx-n" space="4"><mjx-c class="mjx-c34"></mjx-c></mjx-mn><mjx-tstrut></mjx-tstrut></mjx-mtd></mjx-mtr><mjx-mtr><mjx-mtd style="text-align: right; padding-right: 0px; padding-top: 0.15em;"><mjx-mfrac><mjx-frac type="d"><mjx-num><mjx-nstrut type="d"></mjx-nstrut><mjx-mrow><mjx-mn class="mjx-n"><mjx-c class="mjx-c31"></mjx-c><mjx-c class="mjx-c36"></mjx-c></mjx-mn><mjx-mi class="mjx-i"><mjx-c class="mjx-c1D465 TEX-I"></mjx-c></mjx-mi></mjx-mrow></mjx-num><mjx-dbox><mjx-dtable><mjx-line type="d"></mjx-line><mjx-row><mjx-den><mjx-dstrut type="d"></mjx-dstrut><mjx-mi class="mjx-i"><mjx-c class="mjx-c1D466 TEX-I"></mjx-c></mjx-mi></mjx-den></mjx-row></mjx-dtable></mjx-dbox></mjx-frac></mjx-mfrac><mjx-tstrut></mjx-tstrut></mjx-mtd><mjx-mtd style="text-align: left; padding-left: 0px; padding-top: 0.15em;"><mjx-mi class="mjx-n"></mjx-mi><mjx-mo class="mjx-n" space="4"><mjx-c class="mjx-c3D"></mjx-c></mjx-mo><mjx-mi class="mjx-i" space="4"><mjx-c class="mjx-c1D465 TEX-I"></mjx-c></mjx-mi><mjx-mi class="mjx-i"><mjx-c class="mjx-c1D466 TEX-I"></mjx-c></mjx-mi><mjx-tstrut></mjx-tstrut></mjx-mtd></mjx-mtr></mjx-itable></mjx-table></mjx-mtable></mjx-math><mjx-assistive-mml display="block" unselectable="on"><math display="block" xmlns="http://www.w3.org/1998/Math/MathML"><mtable columnalign="right left" columnspacing="0em" displaystyle="true" rowspacing="3pt"><mtr><mtd><mo stretchy="false">|</mo><msup><mi>x</mi><mrow data-mjx-texclass="ORD"><mn>2</mn></mrow></msup><mo>-</mo><msup><mi>y</mi><mrow data-mjx-texclass="ORD"><mn>2</mn></mrow></msup><mo stretchy="false">|</mo></mtd><mtd><mi></mi><mo>=</mo><mn>4</mn></mtd></mtr><mtr><mtd><mfrac><mrow><mn>16</mn><mi>x</mi></mrow><mi>y</mi></mfrac></mtd><mtd><mi></mi><mo>=</mo><mi>x</mi><mi>y</mi></mtd></mtr></mtable></math></mjx-assistive-mml></mjx-container>
$\textbf{(A)}\   2  \qquad 	\textbf{(B)}\   4  \qquad 	\textbf{(C)}\   6  \qquad 	\textbf{(D)}\   8  \qquad 	\textbf{(E)}\   10$


	\item (\textit{AMC 12 12A 2026}) The tree diagram below has nodes labeled with rational numbers. The top node, at level $0$, is $\frac{1}{2}$. Each node $\frac{a}{b}$ sprouts two nodes in the level below it, $\frac{a}{a+b}$ and $\frac{b}{a+b}$. For example, level $1$ has nodes $\frac{1}{1+2} = \frac{1}{3}$ and $\frac{2}{1+2} = \frac{2}{3}$.


 
                    \begin{figure}[H]
                        \centering
                        \includegraphics[width=0.5\textwidth]{Problem/AMC_12_12A_2026_10_Problem_Figure_1.png}
                    \end{figure}
                    


 At what level does the node $\frac{20}{27}$ appear in the diagram?


 

$\textbf{(A)}\   5  \qquad 	\textbf{(B)}\   8  \qquad 	\textbf{(C)}\   12  \qquad 	\textbf{(D)}\   26  \qquad 	\textbf{(E)}\   \text{The node } \dfrac{20}{27} \text{ does not appear in the diagram.}$


	\item (\textit{AMC 12 12A 2026}) An urn contains $6$ red balls, $4$ black balls, and $4$ white balls. Suppose $7$ balls are drawn at random without replacement and a score is computed as follows:


 - $1$ point is given for each red ball up to a maximum of $3$ points.


 - $1$ point is given for each black ball up to a maximum of $2$ points.


 - $1$ point is given for each white ball up to a maximum of $2$ points.


 For example, drawing $1$ red ball, $4$ black balls, and $2$ white balls would earn a score of $1 + 2 + 2 = 5$. What is the probability of getting the lowest possible score?


 $\textbf{(A)}\   \dfrac{1}{858}  \qquad 	\textbf{(B)}\   \dfrac{5}{3432}  \qquad 	\textbf{(C)}\   \dfrac{1}{312}  \qquad 	\textbf{(D)}\   \dfrac{2}{429}  \qquad 	\textbf{(E)}\   \dfrac{5}{264}$


	\item (\textit{AMC 12 12A 2026}) Ethan writes down five congruences:

 \[x \equiv  0 \pmod{2}, x \equiv  0 \pmod{3}, x \equiv 1 \pmod{4}, x \equiv 5 \pmod{6}, x \equiv 7 \pmod{12}.\]
How many positive integers less than or equal to $2026$ satisfy at least one of these five congruences?


 $\textbf{(A)}\   1688  \qquad 	\textbf{(B)}\   1856  \qquad 	\textbf{(C)}\   1857  \qquad 	\textbf{(D)}\   2025  \qquad 	\textbf{(E)}\   2026$


	\item (\textit{AMC 12 12A 2026}) A board has $4$ spaces in a row, which can be filled with numbers. The sides of a fair, $8$-sided die are numbered $1$ through $8$. David rolls the die $4$ times. After each roll, he randomly chooses an empty spot to write the number down on the die. What is the probability that the numbers on the board form a strictly increasing sequence when read from left to right?


 $\textbf{(A)}\   \dfrac{1}{128}  \qquad 	\textbf{(B)}\   \dfrac{5}{432}  \qquad 	\textbf{(C)}\   \dfrac{35}{2048}  \qquad 	\textbf{(D)}\   \dfrac{35}{1024}  \qquad 	\textbf{(E)}\   \dfrac{1}{24}$


	\item (\textit{AMC 12 12A 2026}) Three identical boxes, each with one open face, are made from $1$-inch-thick wood. The outer dimensions of each box are $a$, $b$, and $c$ inches, where $a < b < c$, and each open face appears on a face with outer dimensions $b \times c$.


 
                    \begin{figure}[H]
                        \centering
                        \includegraphics[width=0.5\textwidth]{Problem/AMC_12_12A_2026_14_Problem_Figure_1.png}
                    \end{figure}
                    


 It happens that if the open faces of the two boxes are held together, then the space inside can exactly contain the (rotated) third box, with no room to spare. What is $a + b + c$?


 $\textbf{(A)}\   20  \qquad 	\textbf{(B)}\   22  \qquad 	\textbf{(C)}\   24  \qquad 	\textbf{(D)}\   26  \qquad 	\textbf{(E)}\   28$


	\item (\textit{AMC 12 12A 2026}) Suppose that $\sin a = \frac{1}{4}$ and $\sin b = \frac{1}{8}$ with $0 \le b < a \le  \frac{\pi}{2}$. What is the greatest integer less than $\log_{2}\big( \sin(a+b) \cdot \sin(a-b) \big)$?


 

$\textbf{(A)}\   -5  \qquad 	\textbf{(B)}\   -4  \qquad 	\textbf{(C)}\   0  \qquad 	\textbf{(D)}\   1  \qquad 	\textbf{(E)}\   5$


	\item (\textit{AMC 12 12A 2026}) Gerry is playing the Triples Game. Her task is to partition the $3 \times 3$ grid of squares shown on the left into three sections, each containing three adjoining squares. (Three squares are adjoining if each shares an edge with at least one of the other two.) One point will be awarded for each section containing more $\textsf{X}$'s than $\textsf{O}$'s. Shown on the right is a sample partition that will earn one point in total. A partition that earns two or more points will win the game. What fraction of the possible partitions of this particular grid will result in a win for Gerry?


 
                    \begin{figure}[H]
                        \centering
                        \includegraphics[width=0.5\textwidth]{Problem/AMC_12_12A_2026_16_Problem_Figure_1.png}
                    \end{figure}
                    
$\textbf{(A)}\   \dfrac{3}{10}  \qquad 	\textbf{(B)}\   \dfrac{1}{3}  \qquad 	\textbf{(C)}\   \dfrac{2}{5}  \qquad 	\textbf{(D)}\   \dfrac{4}{9}  \qquad 	\textbf{(E)}\   \dfrac{1}{2}$


	\item (\textit{AMC 12 12A 2026}) An engineer is designing a closely packed field of $7$ wind turbines for electricity production. Each turbine, whose center is marked by a labeled dot in the figure, occupies one of $7$ congruent elliptical regions. Turbines in the same row have centers that are $4$ units apart, and turbines in the same column have centers that are $9$ units apart, as shown in the figure below.


 
                    \begin{figure}[H]
                        \centering
                        \includegraphics[width=0.5\textwidth]{Problem/AMC_12_12A_2026_17_Problem_Figure_1.png}
                    \end{figure}
                    


 The major axis of each ellipse is $m\sqrt{p}$, where $m$ and $p$ are positive integers and $p$ is not divisible by the square of any prime. What is $m + p$?


 $\textbf{(A)}\   6  \qquad 	\textbf{(B)}\   8  \qquad 	\textbf{(C)}\   9  \qquad 	\textbf{(D)}\   10  \qquad 	\textbf{(E)}\   12$


	\item (\textit{AMC 12 12A 2026}) A sphere with center $O$ and radius $9$ cuts three mutually perpendicular planes that meet at point $P$, forming three circular cross-sections on the planes, which have radii $7$, $4$, and $3$. What is $OP$?


 $\textbf{(A)}\   10  \qquad 	\textbf{(B)}\   11  \qquad 	\textbf{(C)}\   12  \qquad 	\textbf{(D)}\   13  \qquad 	\textbf{(E)}\   14$


	\item (\textit{AMC 12 12A 2026}) How many values of $x$ between $0$ and $2\pi$, inclusive, satisfy the equation

 \[\sin x + \sin(2x) + \sin(3x) = 1 + \cos x + \cos(2x)\ ?\]
$\textbf{(A)}\   1  \qquad 	\textbf{(B)}\   2  \qquad 	\textbf{(C)}\   4  \qquad 	\textbf{(D)}\   6  \qquad 	\textbf{(E)}\   8$


	\item (\textit{AMC 12 12A 2026}) How many polynomials $f(x)$ of degree $3$ satisfying

 \[f\left( 3+\sqrt{5} \right) = f\left( 3 - \sqrt{5} \right)\]
are there that have leading coefficient equal to $1$ and all of whose coefficients are integers between $-10$ and $10$, inclusive?


 $\textbf{(A)}\   21  \qquad 	\textbf{(B)}\   42  \qquad 	\textbf{(C)}\   63  \qquad 	\textbf{(D)}\   84  \qquad 	\textbf{(E)}\   105$


	\item (\textit{AMC 12 12A 2026}) Six dogs find a large cookie in the shape of a regular hexagon $ABCDEF$. To divide the cookie, the dogs agree to choose point $P$ uniformly at random in the interior of the hexagon. The dogs will then cut out the segments $\overline{PA}$, $\overline{PB}$, $\overline{PC}$, $\overline{PD}$, $\overline{PE}$, and $\overline{PF}$, dividing the cooking into six triangles, one for each dog. A dog will be happy if and only if its piece has at least $\frac{1}{10}$ of the total area of the hexagon. What is the probability that all six dogs will be happy?


 $\textbf{(A)}\   \dfrac{\sqrt{3}}{50}  \qquad 	\textbf{(B)}\   \dfrac{1}{25}  \qquad 	\textbf{(C)}\   \dfrac{3\sqrt{3}}{50}  \qquad 	\textbf{(D)}\   \dfrac{4}{25}  \qquad 	\textbf{(E)}\   \dfrac{2}{5}$


	\item (\textit{AMC 12 12A 2026}) A flock of $8$ birds is flying in a formation, led by a single leader bird, which has exactly $2$ followers immediately behind. Each non-leader bird follows exactly one bird, and has $0$, $1$, or $2$ followers immediately behind. Two possible formations are shown in the picture below. (When a bird has two followers, swapping its left and right subflocks can produce a different formation, as shown in the figure.)


 

                    \begin{figure}[H]
                        \centering
                        \includegraphics[width=0.5\textwidth]{Problem/AMC_12_12A_2026_22_Problem_Figure_1.png}
                    \end{figure}
                    


 


 What is the total number of possible formations?


 $\textbf{(A)}\   48  \qquad 	\textbf{(B)}\   52  \qquad 	\textbf{(C)}\   76  \qquad 	\textbf{(D)}\   132  \qquad 	\textbf{(E)}\   196$


	\item (\textit{AMC 12 12A 2026}) There is a unique real number $x$ between $0$ and $1$ such that $x^{3} + x^{2} = 1$. For that $x$, there exist unique integers $a$, $b$, and $c$ such that the infinite product

 \[(x+1)(x^{2}+1)(x^{4}+1)(x^{8}+1)(x^{16}+1) \cdots\]
equals $ax^{2} + bx + c$. What is $a+b+c$?


 $\textbf{(A)}\   2  \qquad 	\textbf{(B)}\   3  \qquad 	\textbf{(C)}\   4  \qquad 	\textbf{(D)}\   5  \qquad 	\textbf{(E)}\   6$


	\item (\textit{AMC 12 12A 2026}) Trapezoid $ABCD$ has $AB = 4$, $CD = 5$, and $\angle DAB = \angle ABC = 90^{\circ}$. Circles $\omega_{1}$ and $\omega_{2}$ lie inside the trapezoid as follows: $\omega_{1}$ is tangent to $\overline{AB}$, $\overline{AD}$, and $\overline{DC}$; $\omega_{2}$ is tangent to $\overline{AB}$, $\overline{BC}$, and $\overline{DC}$; and $\omega_{1}$ and $\omega_{2}$ are tangent to each other. The distance between the centers of $\omega_{1}$ and $\omega_{2}$ is $m - n\sqrt{p}$, where $m$, $n$, and $p$ are positive integers such that $p$ is not divisible by the square of any prime. What is $m + n + p$?


 $\textbf{(A)}\   30  \qquad 	\textbf{(B)}\   47  \qquad 	\textbf{(C)}\   53  \qquad 	\textbf{(D)}\   62  \qquad 	\textbf{(E)}\   74$


	\item (\textit{AMC 12 12A 2026}) Let $m$ be a positive integer such that there are at least $50$ positive integers $x$ less than $m$ for which $x^2 - 1$ is a multiple of $m$. What is the sum of the digits of the least possible value of $m$?


 $\textbf{(A)}\   12  \qquad 	\textbf{(B)}\   15  \qquad 	\textbf{(C)}\   16  \qquad 	\textbf{(D)}\   18  \qquad 	\textbf{(E)}\   22$


\end{enumerate}